% ═════════════════════════════════════════════════════════════════════════════
\section{Properties of Functions}
% ═════════════════════════════════════════════════════════════════════════════

These notes examine injectivity, surjectivity, composition, and invertibility
through a detailed worked example involving a function $f : \Rtwo \to \Rtwo$.

% ═════════════════════════════════════════════════════════════════════════════
\section{Definitions}
% ═════════════════════════════════════════════════════════════════════════════

Let $f : A \to B$.

\begin{definition}
$f$ is \defn{injective} (one-to-one) if $f(x_1) = f(x_2) \Rightarrow x_1 = x_2$
for all $x_1, x_2 \in A$.
\end{definition}

\begin{definition}
$f$ is \defn{surjective} (onto) if for every $b \in B$ there exists $a \in A$
with $f(a) = b$. The \defn{range} of $f$ is $\mathrm{Ran}(f) = \{f(a) : a \in A\}$.
$f$ is surjective iff $\mathrm{Ran}(f) = B$.
\end{definition}

\begin{definition}
$f$ is \defn{bijective} if it is both injective and surjective.
\end{definition}

\begin{definition}
$f$ is \defn{invertible} if there exists $g : B \to A$ such that $g \circ f = \mathrm{id}_A$
and $f \circ g = \mathrm{id}_B$. A function is invertible if and only if it is
bijective.
\end{definition}

% ═════════════════════════════════════════════════════════════════════════════
\section{A Worked Example}
% ═════════════════════════════════════════════════════════════════════════════

Let $f : \Rtwo \to \Rtwo$ and $g : \Rtwo \to \Rtwo$ be defined by
\[
    f(x, y) = (y - 2x,\; x^4 + 1),
    \qquad
    g(x, y) = (-y,\; 3 + x).
\]

% ─────────────────────────────────────────────────────────────────────────────
\subsection{Injectivity of \texorpdfstring{$f$}{f}}
% ─────────────────────────────────────────────────────────────────────────────

\begin{proposition}
$f$ is not injective.
\end{proposition}

\begin{proof}
Take $(x_1, y_1) = (1, 10)$ and $(x_2, y_2) = (-1, 6)$.
\begin{align*}
    f(1, 10) &= \bigl(10 - 2(1),\; 1^4 + 1\bigr) = (8, 2), \\
    f(-1, 6) &= \bigl(6 - 2(-1),\; (-1)^4 + 1\bigr) = (8, 2).
\end{align*}
Since $f(1,10) = f(-1,6)$ but $(1,10) \neq (-1,6)$, $f$ is not injective.
\end{proof}

% ─────────────────────────────────────────────────────────────────────────────
\subsection{Surjectivity of \texorpdfstring{$f$}{f} and its range}
% ─────────────────────────────────────────────────────────────────────────────

\begin{proposition}
$f$ is not surjective. Its range is
$\mathrm{Ran}(f) = \{(u, v) \in \Rtwo : v \geq 1\}$.
\end{proposition}

\begin{proof}
We show $(0, 0) \notin \mathrm{Ran}(f)$. Suppose $(x, y) \in \Rtwo$ satisfies
$f(x,y) = (0,0)$. Then
\begin{align}
    y - 2x &= 0, \tag{1}\\
    x^4 + 1 &= 0. \tag{2}
\end{align}
Equation (2) gives $x^4 = -1$, which has no real solution. Hence $(0,0) \notin
\mathrm{Ran}(f)$, and $f$ is not surjective.

For the range: since $x^4 \geq 0$ for all $x \in \R$, we have $x^4 + 1 \geq 1$,
so the second component of $f(x,y)$ is always at least $1$. Conversely, for any
$(u, v) \in \Rtwo$ with $v \geq 1$, set $x = (v-1)^{1/4}$ and $y = u + 2x$;
then $f(x,y) = (y - 2x, x^4 + 1) = (u, v)$. Hence
\[
    \mathrm{Ran}(f) = \{(u, v) \in \Rtwo : v \geq 1\}. \qedhere
\]
\end{proof}

% ─────────────────────────────────────────────────────────────────────────────
\subsection{The composition \texorpdfstring{$f \circ g$}{f∘g}}
% ─────────────────────────────────────────────────────────────────────────────

\begin{proposition}
$(f \circ g)(x, y) = (3 + x + 2y,\; y^4 + 1)$.
\end{proposition}

\begin{proof}
\begin{align*}
    (f \circ g)(x, y)
    &= f\bigl(g(x, y)\bigr) \\
    &= f(-y,\; 3 + x) \\
    &= \bigl((3 + x) - 2(-y),\; (-y)^4 + 1\bigr) \\
    &= (3 + x + 2y,\; y^4 + 1). \qedhere
\end{align*}
\end{proof}

% ─────────────────────────────────────────────────────────────────────────────
\subsection{Invertibility of \texorpdfstring{$f \circ g$}{f∘g}}
% ─────────────────────────────────────────────────────────────────────────────

% CORRECTION: The original argument stated "for f∘g to be invertible, both f
% and g must be bijections". This is false in general: f∘g can be bijective
% even when neither f nor g individually is. The correct argument uses range
% containment: Ran(f∘g) ⊆ Ran(f) ⊊ R², so f∘g cannot be surjective.

\begin{proposition}
$f \circ g$ is not invertible.
\end{proposition}

\begin{proof}
For any function $h : A \to C$ of the form $h = f \circ g$ with $g : A \to B$
and $f : B \to C$, we have $\mathrm{Ran}(f \circ g) \subseteq \mathrm{Ran}(f)$.
Since
\[
    \mathrm{Ran}(f) = \{(u,v) \in \Rtwo : v \geq 1\} \subsetneq \Rtwo,
\]
it follows that $\mathrm{Ran}(f \circ g) \subseteq \mathrm{Ran}(f) \subsetneq
\Rtwo$. Hence $f \circ g$ is not surjective onto $\Rtwo$, and therefore not
bijective, so not invertible.
\end{proof}

